\subsection{Valores Mágicos}

\subsubsection{Regla: enumeradores y constantes en lugar de literales}

\textbf{Regla:} Queda prohibido el uso de valores numéricos o cadenas literales directamente dentro de la lógica del juego. Todo valor con significado en el dominio deberá gestionarse mediante un \texttt{enum} o una constante de clase con nombre descriptivo.

\textbf{Descripción:} Martin identifica los números mágicos como un mal olor de código porque obligan a quien lee el código a memorizar o adivinar el significado de un valor que aparece sin contexto (Martin, 2008).

\textbf{Código correcto:}
\begin{lstlisting}[style=csharpstyle]
public enum TurnState
{
    WaitingForRoll,
    Moving,
    Finished
}

public sealed class TurnManager
{
    private const int MaximumPlayers = 4;

    public bool CanAddPlayer(int currentPlayerCount)
    {
        return currentPlayerCount < MaximumPlayers;
    }
}
\end{lstlisting}

\textbf{Código incorrecto:}
\begin{lstlisting}[style=csharpstyle]
public sealed class TurnManager
{
    public bool CanAddPlayer(int currentPlayerCount)
    {
        return currentPlayerCount < 4;
    }

    public void SetState(int state)
    {
        if (state == 2)
        {
            FinishTurn();
        }
    }

    private void FinishTurn()
    {
    }
}
\end{lstlisting}
